\section{September 28, 2026}

\subsection{Isomorphic vector spaces}

\begin{definition}[Isomorphism]
  An \emph{isomorphism} is an invertible linear map. Two vector spaces
  $V$ and $W$ over the same field $\bb{F}$ are \emph{isomorphic} if
  there exists an isomorphism from $V$ onto $W$. We write $V\cong W$.
\end{definition}

By Proposition~\ref{prop:invertible-iff-bijective}, a linear map is an
isomorphism precisely when it is both injective and surjective.

\begin{proposition}\label{prop:isomorphic-iff-equal-dimension}
  Two finite-dimensional vector spaces over $\bb{F}$ are isomorphic
  if and only if they have the same dimension.
\end{proposition}

\begin{proof}
  Suppose first that $T\colon V\to W$ is an isomorphism. Then $T$ is
  injective and surjective, so $\ker T=\{0\}$ and
  $\operatorname{im}T=W$. By the fundamental theorem of linear maps
  (Theorem~\ref{thm:fundamental-linear-maps}),
  \[
    \dim V=\dim\ker T+\dim\operatorname{im}T
      =0+\dim W=\dim W.
  \]

  Conversely, let $\dim V=\dim W=n$. Choose bases
  $v_1,\ldots,v_n$ of $V$ and $w_1,\ldots,w_n$ of $W$. By the
  linear map lemma (Lemma~\ref{lem:linear-map}), there is a unique
  linear map $T\colon V\to W$ satisfying $T(v_j)=w_j$ for all $j$.
  Explicitly,
  \[
    T(c_1v_1+\cdots+c_nv_n)=c_1w_1+\cdots+c_nw_n.
  \]
  This map is well-defined because coordinates with respect to a
  basis are unique. Since $w_1,\ldots,w_n$ spans $W$, every vector
  in $W$ occurs as an output, so $T$ is surjective. If
  $T(c_1v_1+\cdots+c_nv_n)=0$, then
  \[
    c_1w_1+\cdots+c_nw_n=0.
  \]
  Linear independence of the $w_j$ gives $c_1=\cdots=c_n=0$.
  Thus $\ker T=\{0\}$, and Proposition~\ref{prop:injective-kernel}
  shows that $T$ is injective. Therefore $T$ is an isomorphism.
  When $n=0$, the unique map between the two zero spaces is an
  isomorphism as well.
\end{proof}

\subsection{The space of linear maps and the space of matrices}

Fix ordered bases
\[
  \mathcal B=(v_1,\ldots,v_n)
  \qquad\text{and}\qquad
  \mathcal C=(w_1,\ldots,w_m)
\]
of $V$ and $W$, respectively. To keep track of both bases, write
$[T]_{\mathcal C}^{\mathcal B}$ for the matrix of
$T\in\mathcal L(V,W)$, characterized by
\[
  [T(v)]_{\mathcal C}
    =[T]_{\mathcal C}^{\mathcal B}[v]_{\mathcal B}.
\]
Thus the $j$th column is $[T(v_j)]_{\mathcal C}$, as in
Theorem~\ref{thm:matrix-representation}. The superscript specifies the domain basis and the subscript specifies
the codomain basis. Thus the matrix takes $\mathcal B$-coordinates to
$\mathcal C$-coordinates.

\begin{proposition}\label{prop:matrix-map-isomorphism}
  The map
  \[
    \mathcal M\colon\mathcal L(V,W)\longrightarrow M_{m\times n}(\bb{F}),
    \qquad \mathcal M(T)=[T]_{\mathcal C}^{\mathcal B},
  \]
  is an isomorphism of vector spaces.
\end{proposition}

\begin{proof}
  The calculations in
  Section~\ref{subsec:matrix-operations-linear-maps} show that
  \[
    \mathcal M(T+S)=\mathcal M(T)+\mathcal M(S),
    \qquad
    \mathcal M(\alpha T)=\alpha\mathcal M(T).
  \]
  Hence $\mathcal M$ is linear.

  For injectivity, suppose that $\mathcal M(T)$ is the zero matrix.
  Every column is zero, so for each $j$,
  \[
    T(v_j)=0w_1+\cdots+0w_m=0.
  \]
  Since the $v_j$ form a basis, linearity implies that $T(v)=0$ for
  every $v\in V$. Thus $T$ is the zero map, and
  $\ker\mathcal M=\{0\}$. By
  Proposition~\ref{prop:injective-kernel}, $\mathcal M$ is injective.

  For surjectivity, take an arbitrary matrix
  $A=[A_{ij}]\in M_{m\times n}(\bb{F})$. Prescribe
  \[
    T(v_j)=\sum_{i=1}^m A_{ij}w_i
    \qquad (j=1,\ldots,n).
  \]
  Lemma~\ref{lem:linear-map} gives a linear map $T\in\mathcal L(V,W)$
  with these values, and its matrix is exactly $A$. Therefore
  $\mathcal M$ is surjective. Being linear and bijective, it is an
  isomorphism by Proposition~\ref{prop:invertible-iff-bijective}.
\end{proof}

\begin{corollary}\label{cor:dimension-linear-maps}
  If $V$ and $W$ are finite-dimensional, then $\mathcal L(V,W)$ is
  finite-dimensional and
  \[
    \dim\mathcal L(V,W)=(\dim V)(\dim W).
  \]
\end{corollary}

\begin{proof}
  Set $n=\dim V$ and $m=\dim W$. For $1\leq i\leq m$ and
  $1\leq j\leq n$, let $E_{ij}$ be the matrix with a $1$ in entry
  $(i,j)$ and zeros elsewhere. Every matrix $A=[A_{ij}]$ has the
  unique expansion
  \[
    A=\sum_{i=1}^m\sum_{j=1}^n A_{ij}E_{ij}.
  \]
  Thus the $mn$ matrices $E_{ij}$ form a basis of
  $M_{m\times n}(\bb{F})$. Because $\mathcal M$ is a linear
  isomorphism, the maps $\mathcal M^{-1}(E_{ij})$ form a basis of
  $\mathcal L(V,W)$. Hence this space is finite-dimensional with
  dimension $mn$. If either $m$ or $n$ is zero, the space of linear
  maps is the zero space, and the same formula holds.
\end{proof}

\subsection{Invertible matrices}

\begin{definition}[Invertible matrix]
  A square matrix $A\in M_{n\times n}(\bb{F})$ is \emph{invertible}
  if there exists $B\in M_{n\times n}(\bb{F})$ such that
  \[
    AB=BA=I_n,
  \]
  where $I_n$ is the identity matrix. We call $B$ the \emph{inverse}
  of $A$ and denote it by $A^{-1}$.
\end{definition}

The inverse is unique: if $B$ and $C$ are both inverses of $A$, then
\[
  B=B(AC)=(BA)C=C.
\]

Let $V$ and $W$ be finite-dimensional, and let $T\colon V\to W$
be an isomorphism. By
Proposition~\ref{prop:isomorphic-iff-equal-dimension}, both spaces
have the same dimension $n$. Let
\[
  A=[T]_{\mathcal C}^{\mathcal B},
  \qquad
  B=[T^{-1}]_{\mathcal B}^{\mathcal C}.
\]
By Section~\ref{subsec:composition-matrix-multiplication}, composition
corresponds to matrix multiplication. Thus
$T^{-1}T=\operatorname{Id}_V$ and
$TT^{-1}=\operatorname{Id}_W$ give
\[
  BA=I_n
  \qquad\text{and}\qquad
  AB=I_n.
\]
Consequently,
\[
  [T^{-1}]_{\mathcal B}^{\mathcal C}
    =\bigl([T]_{\mathcal C}^{\mathcal B}\bigr)^{-1}.
\]
Notice that the roles of the domain and codomain bases are reversed
for the inverse map.

\subsection{Change of basis}

Let
\[
  \mathcal B=(v_1,\ldots,v_n)
  \qquad\text{and}\qquad
  \mathcal C=(u_1,\ldots,u_n)
\]
be two ordered bases of the same vector space $V$. Define
\[
  P=[\operatorname{Id}_V]_{\mathcal C}^{\mathcal B},
  \qquad
  Q=[\operatorname{Id}_V]_{\mathcal B}^{\mathcal C}.
\]
Although the identity map leaves each vector unchanged, its matrix
need not be $I_n$ when the domain and codomain use different bases.
Here
\[
  [v]_{\mathcal C}=P[v]_{\mathcal B},
  \qquad
  [v]_{\mathcal B}=Q[v]_{\mathcal C}.
\]
In particular, the columns of $P$ are the old basis vectors
expressed in the new basis:
\[
  P=
  \begin{bmatrix}
    \vert & & \vert \\
    [v_1]_{\mathcal C} & \cdots & [v_n]_{\mathcal C} \\
    \vert & & \vert
  \end{bmatrix}.
\]

\begin{proposition}\label{prop:change-of-basis-inverses}
  The matrices $P$ and $Q$ are invertible, and each is the inverse
  of the other.
\end{proposition}

\begin{proof}
  Changing from $\mathcal B$-coordinates to $\mathcal C$-coordinates
  and then back represents the composition
  $\operatorname{Id}_V\circ\operatorname{Id}_V=\operatorname{Id}_V$
  with $\mathcal B$ used at both ends. Hence $QP=I_n$.
  Reversing the order gives $PQ=I_n$. Thus $Q=P^{-1}$.
\end{proof}

\begin{proposition}[Change-of-basis formula]
  \label{prop:change-of-basis-formula}
  Let $T\in\mathcal L(V,V)$, and let
  \[
    A=[T]_{\mathcal B}^{\mathcal B},
    \qquad
    B=[T]_{\mathcal C}^{\mathcal C}.
  \]
  With $P=[\operatorname{Id}_V]_{\mathcal C}^{\mathcal B}$,
  we have
  \[
    B=PAP^{-1}=PAQ.
  \]
\end{proposition}

\begin{proof}
  Applying $T$ and then changing coordinates has the same effect as
  changing coordinates first and then applying $T$ in the new basis.
  The following diagram expresses this equality:
  \[
    \begin{CD}
      [v]_{\mathcal B} @>{A}>> [T(v)]_{\mathcal B} \\
      @V{P}VV @VV{P}V \\
      [v]_{\mathcal C} @>{B}>> [T(v)]_{\mathcal C}.
    \end{CD}
  \]
  Equivalently, the identity
  $T\circ\operatorname{Id}_V=\operatorname{Id}_V\circ T$ gives
  \[
    BP=PA.
  \]
  Multiplying on the right by $P^{-1}=Q$ yields
  $B=PAP^{-1}=PAQ$.
\end{proof}

To interpret the formula, start with a vector in
$\mathcal C$-coordinates. The rightmost factor $P^{-1}$ converts it
to $\mathcal B$-coordinates, $A$ computes the coordinates of its
image under $T$, and $P$ converts the result back to
$\mathcal C$-coordinates.
